% ============================================================
%  AdaptPlan — PPE Final Report
%  CRMEF Casablanca-Settat | Informatique | 2025-2026
%  Language: English
% ============================================================
\documentclass[12pt,a4paper]{report}

% ---------- Encoding & Language ----------
\usepackage[utf8]{inputenc}
\usepackage[T1]{fontenc}
\usepackage[english]{babel}

% ---------- Page Layout ----------
\usepackage{geometry}
\geometry{a4paper, top=2.3cm, bottom=2.3cm, left=2.8cm, right=2.3cm, headheight=15pt}

% ---------- Typography ----------
\usepackage{setspace}
\setstretch{1.15}

% ---------- Graphics & Color ----------
\usepackage{graphicx}
\usepackage{xcolor}

% ---------- Math ----------
\usepackage{amsmath}
\usepackage{amssymb}

% ---------- Tables & Floats ----------
\usepackage{booktabs}
\usepackage{float}
\usepackage{array}
\usepackage{tabularx}
\usepackage{longtable}

% ---------- Lists ----------
\usepackage{enumitem}

% ---------- Boxes ----------
\usepackage[most]{tcolorbox}

% ---------- Hyperlinks ----------
\usepackage{hyperref}

% ---------- Headers/Footers ----------
\usepackage{fancyhdr}

% ---------- Text utilities ----------
\usepackage{ragged2e}

% ============================================================
%  COLOUR PALETTE
% ============================================================
\definecolor{bleuCRMEF}{HTML}{0B60B0}
\definecolor{bleuFond}{HTML}{EEF5FF}
\definecolor{rougeCRMEF}{HTML}{BE123C}
\definecolor{grisTexte}{HTML}{374151}
\definecolor{vertAccent}{HTML}{0F6E56}
\definecolor{ambreAccent}{HTML}{D97706}
\definecolor{grisLeger}{HTML}{F3F4F6}
\definecolor{bleuSombre}{HTML}{1E3A5F}

% ============================================================
%  HYPERREF SETUP
% ============================================================
\hypersetup{
    colorlinks  = true,
    linkcolor   = bleuCRMEF,
    urlcolor    = bleuCRMEF,
    citecolor   = bleuCRMEF,
    pdftitle    = {[YOUR PROJECT TITLE HERE: The Design of...]},
    pdfauthor   = {[CANDIDATE NAME(S) HERE]},
    pdfsubject  = {Designing AI-Powered Applications for Supporting Learners with Special Needs During Examinations},
}

% ============================================================
%  HEADERS & FOOTERS
% ============================================================
\pagestyle{fancy}
\fancyhf{}
\fancyhead[L]{\small\color{grisTexte} Supervised Personal Project (PPE) — AdaptPlan}
\fancyhead[R]{\small\color{grisTexte} CRMEF Casablanca-Settat}
\fancyfoot[C]{\thepage}
\renewcommand{\headrulewidth}{0.4pt}
\renewcommand{\footrulewidth}{0.4pt}

% ============================================================
%  CHAPTER TITLE STYLE
% ============================================================
\usepackage{titlesec}
\titleformat{\chapter}[display]
  {\normalfont\Large\bfseries\color{bleuCRMEF}}
  {\chaptertitlename\ \thechapter}{10pt}{\Huge}
\titlespacing*{\chapter}{0pt}{20pt}{30pt}

% ============================================================
%  CUSTOM ENVIRONMENTS — Classic Academic Style (No AI Border Colors)
% ============================================================
\newenvironment{infobox}[1][]{%
  \begin{quote}\small\itshape\color{grisTexte}%
}{%
  \end{quote}%
}

\newenvironment{alertbox}[1][]{%
  \begin{quote}\small\bfseries\color{grisTexte}%
}{%
  \end{quote}%
}

% ============================================================
%  DOCUMENT START
% ============================================================
\begin{document}

% ============================================================
%  SECTION 1 — COVER PAGE
% ============================================================
\begin{titlepage}
  \thispagestyle{empty}

  \noindent
  \begin{minipage}[b]{0.48\linewidth}
    \includegraphics[height=2.2cm]{CRMEF-LOGO.png}
  \end{minipage}%
  \hfill
  \begin{minipage}[b]{0.48\linewidth}
    \raggedleft
    \includegraphics[height=2.2cm]{ma-logo.jpeg}
  \end{minipage}

  \vspace{1mm}
  \noindent\rule{\linewidth}{1.2pt}
  \vspace{3mm}

  \vspace{4mm}

  \begin{center}
    \setstretch{1.3}
    \fontsize{12}{16}\selectfont
    Teacher Training Cycle \\
    Training Track for Lower Secondary Education Teachers \\
    \textbf{Speciality: [YOUR SPECIALITY HERE, e.g. Computer Science]}
  \end{center}

  \vspace{6mm}

  \begin{center}
    \fontsize{15}{18}\selectfont\bfseries\color{bleuCRMEF}
      Supervised Personal Project (PPE)}
  \end{center}

  \vspace{4mm}

  \begin{center}
    \noindent\rule{0.9\linewidth}{1pt} \\[10pt]
    {\fontsize{16}{20}\selectfont\bfseries\color{bleuCRMEF}
      \setstretch{1.2}
      \begin{tabular}{c}
        [YOUR PROJECT TITLE HERE] \\
        [Sub-title / Explanatory Technical Tagline]
      \end{tabular}
    } \\[10pt]
    \noindent\rule{0.9\linewidth}{1pt}
  \end{center}

  \vspace{8mm}

  \begin{center}
    \fontsize{11}{14}\selectfont
    \begin{tabular}{p{0.45\linewidth} p{0.45\linewidth}}
      \textbf{Prepared by:} & \textbf{Supervised by:} \\
      [Candidate Name(s)] & Prof. [Supervisor Name]
    \end{tabular}
  \end{center}

  \vspace{8mm}

  \begin{center}
    \fontsize{11}{14}\selectfont
    \begin{tabular}{p{0.9\linewidth}}
      \hline
      \textbf{Examination Committee:} \\
      \\[6pt] 1.\ \dotfill \\
      \\[4pt] 2.\ \dotfill \\
      \\[4pt] 3.\ \dotfill \\
      \hline
    \end{tabular}
  \end{center}

  \vfill

  \begin{center}
    {\large\bfseries\color{bleuCRMEF}
      Training Year: [2025--2026]}
  \end{center}
\end{titlepage}

% ============================================================
%  ACKNOWLEDGEMENTS
% ============================================================
\newpage
\pagenumbering{roman}
\phantomsection
\addcontentsline{toc}{section}{Acknowledgements}

\vspace*{1.5cm}
\begin{center}
  {\LARGE\bfseries\color{bleuCRMEF} Acknowledgements}
\end{center}
\vspace{1.0cm}

{\large
\begin{spacing}{1.35}
\justifying\color{grisTexte}

First and foremost, I praise and thank Almighty Allah for granting me the strength, health, and wisdom to successfully complete this project and my qualifying training.

I extend my profound gratitude to my beloved family for their endless sacrifices, patience, and prayers, which have been my cornerstone. To my dedicated friends and colleagues, I express my sincere appreciation for their constant encouragement throughout this journey.

I extend my deepest thanks to my supervisor, \textbf{Prof. BROUMI Said}, for his invaluable guidance, constructive feedback, and sustained support throughout the entire project cycle. His academic insights enabled me to sharpen the project’s objectives and continuously elevate the quality of the work produced.

I would also like to express my sincere recognition to \textbf{Mr. El Mostapha Atoubi}, Head of the Branch, for his dedicated accompaniment throughout this training program, his constant availability, and his unwavering commitment to the success of the teacher-trainees.

I warmly thank the administration, faculty, and the entire pedagogical team of the \textbf{Centre Régional des Métiers de l'Éducation et de la Formation (CRMEF) — Casablanca-Settat}, for the exceptional richness of the teachings shared and the professional contributions that enriched our qualifying training.

Finally, I extend my gratitude to the management of our placement school, \textbf{Al Mokhtar Soussi Middle School}, for their warm reception and the favorable working conditions provided during the practicum. Special thanks go to my classroom mentor, \textbf{Prof. IBISSOUR Rachid}, whose practical advice, field guidance, and pedagogical insights greatly informed my professional reflection as a future educator.

\end{spacing}
}

\newpage

% ============================================================
%  SECTION 2 — TABLE OF CONTENTS
% ============================================================
\tableofcontents
\newpage

% ============================================================
%  SECTION 3 — GENERAL INTRODUCTION
% ============================================================
\pagenumbering{arabic}
\setcounter{page}{1}

\chapter{General Introduction}

\section{Introduction and Background}
Provide an overview of the educational system, the strategic roadmap, or the national reform directives that frame your Supervised Personal Project (PPE). Detail the broad significance of inclusive education and digital integration in your speciality.

\section{Problem Statement}
Clearly state the pedagogical, administrative, or technical problem that this project aims to solve. Explain the current limitations, bottlenecks, or obstacles encountered in practice (e.g., in classrooms or placement schools).

\section{Project Objectives}
Outline the primary goals and specific operational targets of the proposed system:
\begin{itemize}
  \item \textbf{Operational Goal 1:} Describe the main system deliverable.
  \item \textbf{Operational Goal 2:} Detail any analytical, predictive, or classification outputs.
  \item \textbf{Operational Goal 3:} Mention environmental or accessibility accommodations.
\end{itemize}

\section{Structure of the Report}
Briefly summarize the contents of each subsequent chapter of this report to guide the reader.

% ============================================================
\chapter{Theoretical Framework & Methodology}
% ============================================================
\section{Pedagogical Theories and Directives}
Provide a review of the theoretical models and educational frameworks underpinning your project.

\subsection{Theory Name (e.g., Cognitive Load Theory)}
Discuss how the human working memory processes information within this context. Address:
\begin{itemize}
  \item \textbf{Intrinsic Load:} Conceptual complexity of the tasks.
  \item \textbf{Extraneous Load:} Formatting, instruction organization, and presentation barriers.
  \item \textbf{Germane Load:} Schema construction processes.
\end{itemize}

\subsection{Algorithmic Guardrails for Construct Preservation}
Detail how your software implementation operationally enforces these theoretical principles (e.g., via prompt constraints, validation wrappers, or interface controls).

\subsection{Zone of Proximal Development & Scaffolding (Vygotsky)}
Describe how your system implements progressive scaffolding or cues to measure dynamic potential.

\subsection{Universal Design for Learning (UDL)}
Address UDL's three core principles: multiple means of representation, multiple means of action/expression, and multiple means of engagement.

\section{Regulatory and Institutional Framework}
Discuss national ministerial circulars, frameworks, or guidelines governing your project area.

% ============================================================
\chapter{Methodology and Action Plan}
% ============================================================
\section{Target Population and Needs Assessment}
Detail the target learners, teachers, or classrooms participating in the study.

\section{Methodological Workflow}
Explain the design science research methodology, data collection, or development workflow.
\begin{figure}[H]
  \centering
  % \includegraphics[width=0.8\textwidth]{placeholder_workflow.png}
  \caption{Methodological Workflow and Action Plan (Placeholder)}
  \label{fig:workflow}
\end{figure}

\section{Dataset Acquisition and Feature Definition}
Detail the input features, categorical variables, target outputs, and balanced class distributions:
\begin{table}[H]
  \centering
  \small
  \caption{Input and Output Feature Specifications}
  \label{tab:features}
  \begin{tabular}{lll}
    \toprule
    \textbf{Category} & \textbf{Feature Variable} & \textbf{Type / Range} \\
    \midrule
    Inputs & \texttt{feature\_1} & Categorical \\
    Inputs & \texttt{feature\_2} & Numeric \\
    Outputs & \texttt{target\_1} & Binary \\
    \bottomrule
  \end{tabular}
\end{table}

% ============================================================
\chapter{Factual Results}
% ============================================================
\section{System Architecture & Class Structure}
Explain the UML models, database schemas, and logic flow of your modules.

\section{Quantitative Model Performance}
Present your testing runs, hyperparameter tuning curves, and cross-validation metrics.
\begin{figure}[H]
  \centering
  % \includegraphics[width=0.8\textwidth]{plot_cv_curve.png}
  \caption{Model Cross-Validation Tuning Curve (Placeholder)}
  \label{fig:cv_curve}
\end{figure}

On the independent test set, the final model achieved the key metrics presented in Table~\ref{tab:model_metrics}.

\begin{table}[H]
  \centering
  \small
  \caption{Comprehensive Cross-Validation and Target Label Performance Metrics}
  \label{tab:model_metrics}
  \vspace{0.3cm}
  \begin{tabular}{lc}
    \toprule
    \textbf{Evaluation Metric (5-Fold Shuffled CV)} & \textbf{Achieved Accuracy Score} \\
    \midrule
    \textbf{Mean Exact Match (Subset) Accuracy} & \textbf{97.9594\%} \\
    \midrule
    \textit{Per-Label Target Classifier Accuracies:} & \\
    \ \ -- \texttt{target\_label\_1} & 100.0000\% \\
    \ \ -- \texttt{target\_label\_2} & 99.4708\% \\
    \bottomrule
  \end{tabular}
\end{table}

% ============================================================
\chapter{Analysis and Interpretation of Results}
% ============================================================
\section{Statistical Interpretation and Model Robustness}
Analyze the predictive stability of your system. Clarify strict metrics versus per-label values.

\begin{infobox}[title={Understanding Exact Match vs. Per-Label Accuracy}]
  Use this infobox to explain any mathematical nuances, data-ordering artifacts, or performance differences between shuffled and non-shuffled cross-validation.
\end{infobox}

\section{Pedagogical Validation of the System Outputs}
Validate that the correlation matrices and recommendations align with the theoretical frameworks and expert pedagogical rules.

\section{Interpretation of System Scaling}
Provide statistical evidence of differentiation and progressive scaffolding relative to severity or difficulty categories.

% ============================================================
\chapter{Discussion of Results}
% ============================================================
\section{Achieved Objectives vs. Initial Objectives}
Evaluate the final platform deliverables against the initial Stage 3 specifications.

\section{Engineering Challenges and Resolutions}
Document the primary development and integration constraints encountered and how they were resolved.

\subsection{Database Performance and Concurrency}
Explain query optimizations, asynchronous tasks, or concurrency locks.

\subsection{Security and Data Privacy}
Detail data scrubbers, anonymization algorithms, or local parsing engines.

\subsection{Exception Handling and System Resilience}
Discuss defensive string coercion, graceful catch-all fallbacks, and handling of edge-case inputs.

% ============================================================
\chapter{Conclusions and Recommendations}
% ============================================================
\section{Synthesis of Achievements}
Provide a summary of the project's practical, pedagogical, and technical achievements.

\section{Limitations and Future Perspectives}
Discuss future roadmaps, LMS integrations, or compiling to client-side runtimes.

% ============================================================
\chapter{Retrospective: Summary and Professional Lessons Learned}
% ============================================================
\section{Practicum Summary and Classroom Context}
Describe the placement school, your classes, and how the project directly connected with your practical teaching.

\section{Development of Professional Competencies}
Reflect on how this project has shaped your teaching practice, classroom management, and professional reflection as an educator.

% ============================================================
\chapter{References}
% ============================================================
\begin{thebibliography}{99}
  \bibitem{unesco2009}
  UNESCO (2009). \textit{Policy Guidelines on Inclusion in Education}. Paris: United Nations Educational, Scientific and Cultural Organization.

  \bibitem{cast2024}
  CAST (2024). \textit{Universal Design for Learning Guidelines version 3.0}. Wakefield, MA: Center for Applied Special Technology.
\end{thebibliography}

% ============================================================
\chapter*{Appendices}
\addcontentsline{toc}{chapter}{Appendices}
% ============================================================
Include supplemental materials, system screenshots, user guides, or database SQL dumps here.

\end{document}
